% 第 11 章 投资建议与组合构建

\section{投资建议：标配偏积极，江波龙上调至买入}

结合 ch08 的完整估值更新，AStock 当前组合结论是\textbf{标配偏积极 AI 存储链 A 股映射}：\textbf{江波龙中报业绩预告（上半年净利92--110亿，同比暴增622--744倍）验证存储涨价周期盈利弹性}，北京君正机构调研确认存储芯片持续涨价、客户接受度高。覆盖标的权重加权基准空间从-17.0\%大幅提升至\textbf{+8.5\%}，江波龙上调至"买入"评级，北京君正上调至"增持"。

\begin{riskbox}[组合建议]
\textbf{当前动作：标配偏积极 / 江波龙买入。} \textbf{核心推荐江波龙（存储涨价最直接受益者，目标价850--950元）}，重点关注北京君正（存储涨价确认+车载存储龙头）。可保留澜起科技、北方华创作为中性观察；对中微公司、拓荆科技、沪硅产业、长电科技、南大光电降低权重。兆易创新Q1净利+523\%，\textbf{等待中报确认上修}，暂维持中性。
\end{riskbox}

\section{组合分层与行动规则}

本报告的组合建议不是主题仓位表，而是基于目标价空间和证据质量的行动框架。评级含义如下：买入 = 基准空间 \(\geq\)20\%；增持 = 8--20\%；中性 = -15--8\%；减持 = \(-15\%\) 以下；观察 = 无法形成当前目标价。

\begin{exhibitbox}[表 11-1 \quad 投资行动清单 / Investment Action List (2026-07-03更新)]
\centering
\scriptsize
\renewcommand{\arraystretch}{1.22}
\begin{tabularx}{\exhibitboxwidth}{>{\bfseries\raggedright\arraybackslash}p{1.5cm} >{\centering\arraybackslash}p{1.5cm} >{\raggedleft\arraybackslash}p{0.86cm} >{\raggedleft\arraybackslash}p{0.86cm} >{\raggedleft\arraybackslash}p{0.95cm} >{\centering\arraybackslash}p{0.85cm} >{\raggedright\arraybackslash}X}
\toprule
\textbf{标的} & \textbf{建议权重} & \textbf{现价} & \textbf{目标价} & \textbf{空间} & \textbf{评级} & \textbf{组合动作} \\
\midrule
\rowcolor{riskgreen!12}
\textbf{江波龙} & \textbf{核心 18--20\%} & \textbf{618.02} & \textbf{900.00} & \textbf{+45.6\%} & \stance{riskgreen}{\textbf{买入}} & \textbf{\astockstar{} 中报预告92--110亿验证涨价弹性，目标价850--950} \\
\rowcolor{riskgreen!6}
北京君正 & 卫星 8--10\% & 259.56 & 280.00 & +7.9\% & \stance{riskgreen}{增持} & 存储涨价确认，车载+工业，逢低增持 \\
澜起科技 & 核心 16--18\% & 266.80 & 232.05 & -13.0\% & \stance{riskamber}{中性} & 不追涨；PCIe6.x/CXL3.x测试中，等待CXL收入桥 \\
北方华创 & 核心 13--15\% & 816.00 & 869.40 & +6.5\% & \stance{riskamber}{中性} & 设备平台龙头，订单确定性最高，若订单上修可升至增持 \\
中微公司 & 低配 6--8\% & 413.69 & 277.20 & -33.0\% & \stance{riskred}{减持} & 14.7亿参设产业基金，降低权重，等待PE消化 \\
拓荆科技 & 低配 5--7\% & 832.00 & 592.20 & -28.8\% & \stance{riskred}{减持} & \textbf{停牌}筹划收购无锡尚积，复牌后跟踪进展 \\
沪硅产业 & 低配 3--5\% & 34.53 & 24.78 & -28.2\% & \stance{riskred}{减持} & 百亿加码硅片，PE屏蔽，PS/PB显示过高 \\
深科技 & 观察 4--6\% & 55.91 & 45.50 & -18.6\% & \stance{riskamber}{中性} & HBM尚在研发阶段，维持低权重观察 \\
长电科技 & 低配 3--5\% & 90.88 & 55.60 & -38.8\% & \stance{riskred}{减持} & 78亿临港封测项目，AI封装收入兑现前不加仓 \\
通富微电 & 观察 3--4\% & 64.80 & 57.60 & -11.1\% & \stance{riskamber}{中性} & 42.2亿定增获批，AMD绑定需转化为利润率证据 \\
\rowcolor{riskamber!6}
兆易创新 & 观察 3--4\% & 677.77 & 505.80 & -25.4\% & \stance{riskamber}{中性} & Q1净利+523\%，\textbf{等待中报确认上修}，若预告大超预期可上调 \\
南大光电 & 低配 1--2\% & 80.18 & 51.35 & -36.0\% & \stance{riskred}{减持} & 材料弹性需订单而非概念验证 \\
\midrule
\rowcolor{panelgrey}
\multicolumn{4}{>{\cellcolor{panelgrey}\bfseries\raggedright}p{5.4cm}}{\textbf{组合层面}} & \textbf{+8.5\%} & \textbf{标配偏积极} & \textbf{江波龙上调至买入，存储涨价周期确认} \\
\bottomrule
\end{tabularx}
\end{exhibitbox}
\sourcenote{\textbf{2026-07-03更新}：建议权重为 AStock 内部研究组合框架，不是交易指令。江波龙目标价基于中报业绩预告（92--110亿）上修，详见第八章。北京君正新增为卫星层重点覆盖。}

\section{上修与下修触发}

\begin{exhibitbox}[表 11-2 \quad 估值触发条件 / Re-rating Triggers]
\centering
\scriptsize
\renewcommand{\arraystretch}{1.22}
\begin{tabularx}{\exhibitboxwidth}{>{\bfseries\raggedright\arraybackslash}p{1.55cm} >{\raggedright\arraybackslash}X >{\raggedright\arraybackslash}X >{\raggedright\arraybackslash}X}
\toprule
\textbf{标的组} & \textbf{上修触发} & \textbf{维持中性条件} & \textbf{下修/退出触发} \\
\midrule
\rowcolor{riskgreen!6}
\textbf{江波龙/北京君正} & \textbf{存储涨价持续（合约价QoQ>+10\%），中报/三季报继续大超预期} & 涨价兑现但毛利率环比增速放缓 & 存储价格提前回落，或库存减值风险 \\
澜起/北方 & CXL 或设备订单超预期，2027E EPS 上修 15\% 以上 & EPS 按当前一致预期兑现，估值温和消化 & CXL 放量延后、订单确认低于预期或 PE 继续压缩 \\
中微/拓荆 & 海外客户突破或国产设备份额上修，毛利率稳定 & 订单兑现但估值仍高于海外可比两倍以上 & 订单延后、毛利率下滑或客户验证失败 \\
封测/模组 & AI 封装收入占比超过 10\%，利润率持续改善 & 收入增长但毛利率未明显抬升 & AI 订单低于预期或价格竞争压缩利润率 \\
兆易创新 & \textbf{中报预告大超预期（净利>25亿），利基DRAM涨价持续} & Q1高增长但Q2环比放缓 & 存储价格回落，或毛利率下滑 \\
存储/材料 & 利基 DRAM/NOR 涨价持续超两季，材料收入进入公司披露 & 涨价兑现但 PE 仍高于成长匹配区间 & 存储价格提前回落，材料验证失败或收入贡献无法量化 \\
\bottomrule
\end{tabularx}
\end{exhibitbox}
\sourcenote{上修/下修触发条件基于 AStock 盈利预测模型与历史估值区间。江波龙/北京君正上调评级需存储涨价持续（合约价 QoQ >+10\%）与中报/三季报业绩验证。}

\section{最终结论}

\textbf{当前结论：标配偏积极，江波龙上调至买入。} 本报告已完成2026-07-03收盘价基础上的完整估值更新。\textbf{核心变化}：江波龙发布中报业绩预告（上半年净利92--110亿，同比暴增622--744倍），验证存储涨价周期盈利弹性，目标价从660元上修至850--950元，评级从"中性"上调至"买入"。北京君正机构调研确认存储芯片持续涨价、客户接受度高，上调至"增持"。

AI 存储链仍是必须跟踪的产业方向，\textbf{存储涨价主线已获业绩验证}。组合上建议\textbf{标配偏积极}，核心配置江波龙（存储涨价最直接受益者），重点关注北京君正（车载存储+涨价确认），同时保留北方华创、澜起科技等产业主线质量最高的覆盖标的。\par

\vspace{0.5em}
\noindent\textbf{后续跟踪重点}：
\begin{enumerate}[label=\protect\gmark]
  \item 7月中下旬：兆易创新、北京君正中报业绩预告
  \item 7月下旬-8月上旬：中报业绩预告密集披露期
  \item 持续跟踪：DRAM/NAND Flash合约价走势（DRAMeXchange, TrendForce）
  \item 拓荆科技复牌后：收购无锡尚积进展
\end{enumerate}
